\documentclass[11pt]{article}
% ===== Shared preamble for Calendar notes =====
\usepackage[margin=1in]{geometry}
\usepackage{amsmath,amssymb}
\usepackage{enumitem}
\usepackage{xcolor}
\usepackage{tcolorbox}
\tcbuselibrary{skins,breakable}
\usepackage{titlesec}
\usepackage{fancyhdr}
\usepackage{hyperref}
\usepackage{array}
\usepackage{booktabs}
\usepackage{longtable}
\usepackage{pifont} % for \ding

% ---- Colors ----
\definecolor{primary}{HTML}{1B3A6B}
\definecolor{accent}{HTML}{2E7D32}
\definecolor{warnclr}{HTML}{B45309}
\definecolor{tipclr}{HTML}{0F5D8C}
\definecolor{lightbg}{HTML}{F4F6F8}

% ---- Section styling ----
\titleformat{\section}
  {\normalfont\Large\bfseries\color{primary}}
  {\thesection}{0.5em}{}[{\color{primary}\titlerule[1pt]}]
\titleformat{\subsection}
  {\normalfont\large\bfseries\color{primary}}
  {\thesubsection}{0.5em}{}

% ---- Header/Footer ----
\pagestyle{fancy}
\fancyhf{}
\renewcommand{\headrulewidth}{0.4pt}
\fancyfoot[C]{\thepage}

% ---- Hyperref ----
\hypersetup{
  colorlinks=true,
  linkcolor=primary,
  urlcolor=tipclr,
  citecolor=accent
}

% ---- Custom boxes ----
% Formula / Key concept box
\newtcolorbox{formulabox}[1][]{
  enhanced, breakable,
  colback=lightbg, colframe=primary,
  boxrule=0.8pt, arc=2pt,
  left=8pt, right=8pt, top=6pt, bottom=6pt,
  title={\textbf{#1}},
  fonttitle=\bfseries\color{white},
  colbacktitle=primary,
  attach boxed title to top left={xshift=8pt, yshift=-2mm},
  boxed title style={boxrule=0pt, arc=2pt}
}

% Tip / trick box
\newtcolorbox{tipbox}[1][Tip]{
  enhanced, breakable,
  colback=blue!4, colframe=tipclr,
  boxrule=0.8pt, arc=2pt,
  left=8pt, right=8pt, top=6pt, bottom=6pt,
  title={\textbf{\textcolor{white}{\ding{72}\ #1}}},
  fonttitle=\bfseries, colbacktitle=tipclr,
}

% Warning / common mistake box
\newtcolorbox{warnbox}[1][Watch Out]{
  enhanced, breakable,
  colback=orange!6, colframe=warnclr,
  boxrule=0.8pt, arc=2pt,
  left=8pt, right=8pt, top=6pt, bottom=6pt,
  title={\textbf{\textcolor{white}{#1}}},
  fonttitle=\bfseries, colbacktitle=warnclr,
}

% Example / worked question box
\newtcolorbox{examplebox}[1][Example]{
  enhanced, breakable,
  colback=green!4, colframe=accent,
  boxrule=0.8pt, arc=2pt,
  left=8pt, right=8pt, top=6pt, bottom=6pt,
  title={\textbf{\textcolor{white}{#1}}},
  fonttitle=\bfseries, colbacktitle=accent,
}


\title{\color{primary}\textbf{Calendar --- Practice Questions}\\[4pt]\large Medium to Difficult, With Full Solutions}
\author{Aptitude Grind}
\date{}

\lhead{\small Calendar --- Practice Questions}
\rhead{\small Aptitude Grind}

\begin{document}
\maketitle
\thispagestyle{fancy}

\noindent\textit{Difficulty is marked as \textbf{[M]} Medium or \textbf{[H]} Hard. Solve on your own
first, then check the solution. Solutions use the odd-days method from \texttt{01-fundamentals.tex}.}

\bigskip
\hrule
\bigskip

\subsection*{Q1 [M] --- Direct Date Lookup}
What was the day of the week on 26 January 1950?

\begin{tipbox}[Solution]
Reference: 1 January 1900 was a \textbf{Monday}.
\begin{itemize}[leftmargin=1.4em]
  \item Years 1900--1949 (50 years elapsed up to start of 1950... but we need up to 1 Jan 1950): count
    odd days for years 1900 to 1949 inclusive of leap years.
  \item Leap years 1900--1949: 1904, 1908, \ldots, 1948 --- but 1900 is \emph{not} leap (century rule)
    $\Rightarrow$ 12 leap years, 38 ordinary years.
  \item Odd days $= (12 \times 2) + (38 \times 1) = 24 + 38 = 62 \equiv 62 - 56 = 6 \pmod 7$
  \item Monday $+ 6$ odd days $\to$ \textbf{Sunday} (this gives the day for 1 Jan 1950)
  \item Add 25 more days to reach 26 Jan: $25 \bmod 7 = 4$ odd days
  \item Sunday $+ 4 \to$ \textbf{Thursday}
\end{itemize}
\textbf{Answer: Thursday}
\end{tipbox}

\subsection*{Q2 [M] --- N Days After}
Today is Wednesday. What day will it be after 90 days?

\begin{tipbox}[Solution]
$90 \bmod 7 = 90 - 84 = 6$ odd days.
Wednesday $+ 6$ days $\to$ \textbf{Tuesday}.
\textbf{Answer: Tuesday}
\end{tipbox}

\subsection*{Q3 [M] --- Same Date, Different Year}
If 15 March 2015 was a Sunday, what day was 15 March 2016?

\begin{tipbox}[Solution]
2015 is an ordinary year, but the gap 15 Mar 2015 $\to$ 15 Mar 2016 \textbf{includes 29 Feb 2016}
(since 2016 is a leap year and the leap day falls before 15 March). So this interval has
\textbf{2 odd days}, not 1.
Sunday $+ 2 \to$ \textbf{Tuesday}.
\textbf{Answer: Tuesday}
\end{tipbox}

\begin{warnbox}[Why This Trips People Up]
The "add 1 odd day per year" shortcut only works when the year in question is ordinary \textbf{and}
the leap day (if any) doesn't fall between the two matching dates. Always check whether 29 February
lies inside the specific date range you're bridging --- not just whether the calendar year is leap.
\end{warnbox}

\subsection*{Q4 [M] --- Same Calendar Year}
The calendar for the year 2023 will be the same as the calendar for which future year?

\begin{tipbox}[Solution]
2023 is ordinary. Track running odd days year by year starting the count from 2023$\to$2024:
\begin{center}
\begin{tabular}{lcc}
\toprule
Transition & Odd days added & Running total (mod 7) \\
\midrule
2023 $\to$ 2024 (2023 ordinary) & 1 & 1 \\
2024 $\to$ 2025 (2024 leap) & 2 & 3 \\
2025 $\to$ 2026 (2025 ordinary) & 1 & 4 \\
2026 $\to$ 2027 (2026 ordinary) & 1 & 5 \\
2027 $\to$ 2028 (2027 ordinary) & 1 & 6 \\
2028 $\to$ 2029 (2028 leap) & 2 & 1 \\
2029 $\to$ 2030 (2029 ordinary) & 1 & 2 \\
2030 $\to$ 2031 (2030 ordinary) & 1 & 3 \\
2031 $\to$ 2032 (2031 ordinary) & 1 & 4 \\
2032 $\to$ 2033 (2032 leap) & 2 & 6 \\
2033 $\to$ 2034 (2033 ordinary) & 1 & 0 \\
\bottomrule
\end{tabular}
\end{center}
Running total first hits a multiple of 7 at \textbf{2034}, and 2034 is also an ordinary year
(matching 2023's type). \textbf{Answer: 2034}
\end{tipbox}

\subsection*{Q5 [H] --- Century Reasoning}
Prove that the last day of a century (year ending in 00) can never be a Tuesday, Thursday, or Saturday.

\begin{tipbox}[Solution]
The last day of the $n$-th century corresponds to a cumulative odd-day count from a fixed starting
point equal to $5n \bmod 7$ for $n$ not a multiple of 4, and $0$ when $n$ is a multiple of 4
(from the 400-year rule). Evaluate $5n \bmod 7$ for $n = 1, 2, 3$ (since $n=4$ gives 0):
\[
n=1: 5 \quad n=2: 10 \equiv 3 \quad n=3: 15 \equiv 1 \quad n=4: 0
\]
So the only possible odd-day residues for a century's last day are $\{0, 1, 3, 5\}$, which --- using
the reference that year 1's last day (31 Dec, year 0 effectively) corresponds to Sunday being residue
0 --- map to a fixed subset of just four weekdays. Direct computation (as shown in standard references)
confirms centuries can only end on \textbf{Sunday, Monday, Wednesday, or Friday}, never Tuesday,
Thursday, or Saturday.
\textbf{Answer: Because only 4 distinct odd-day residues (0, 1, 3, 5) are possible for century
boundaries, only 4 of the 7 weekdays are reachable --- and Tuesday, Thursday, Saturday are never
among them.}
\end{tipbox}

\subsection*{Q6 [H] --- Reverse Problem}
1 January of a certain year is a Friday. 1 January of the following year is a Sunday. What can you
conclude about the certain year, and what day is 1 January two years later?

\begin{tipbox}[Solution]
Going from Friday to Sunday is a shift of \textbf{2 odd days} $\Rightarrow$ the certain year must be a
\textbf{leap year} (only leap years contribute 2 odd days). So the certain year has 366 days, and the
following year is ordinary, contributing 1 odd day.
Sunday $+ 1 \to$ \textbf{Monday} for 1 January two years later.
\textbf{Answer: The certain year is a leap year; 1 January two years later is a Monday.}
\end{tipbox}

\subsection*{Q7 [H] --- Multi-Step Combined Problem}
If 1 January 2001 was a Monday, find the day of the week on 29 February 2004, and verify 2004 is
indeed a leap year using the divisibility rule.

\begin{tipbox}[Solution]
\textbf{Leap year check:} $2004 / 4 = 501$ exactly, and 2004 is not a century year, so \textbf{2004 is
a leap year}. \checkmark

\textbf{Odd days from 1 Jan 2001 to 1 Jan 2004} (years 2001, 2002, 2003 --- all ordinary):
$3 \times 1 = 3$ odd days. Monday $+ 3 \to$ \textbf{Thursday} (this is 1 Jan 2004).

\textbf{Odd days from 1 Jan 2004 to 29 Feb 2004:}
January has 31 days $\to 31 \bmod 7 = 3$ odd days (this brings us to 1 Feb).
Add 28 more days to reach 29 Feb: $28 \bmod 7 = 0$ odd days.
Total: $3 + 0 = 3$ odd days.

Thursday $+ 3 \to$ \textbf{Sunday}.
\textbf{Answer: 29 February 2004 was a Sunday.}
\end{tipbox}

\subsection*{Q8 [H] --- Century-Spanning Problem}
If 1 January 2001 was a Monday, what day of the week is 1 January 2101?

\begin{tipbox}[Solution]
We need the total odd days contributed by the years 2001 through 2100 (100 full years elapsing).

\textbf{Leap years in 2001--2100:} every year divisible by 4 is a leap year \emph{unless} it's a
century year not divisible by 400. In this range: 2004, 2008, \ldots, 2096 are leap
$\Rightarrow \dfrac{2096 - 2004}{4} + 1 = 24$ leap years. The only century year in range, 2100, is
\textbf{not} divisible by 400, so it is \emph{not} leap (and it's already excluded from the count
above since it isn't divisible by 4 in the "leap" sense we need to re-verify --- 2100 is divisible by
4 but excluded by the century rule, so it contributes as an \textbf{ordinary} year instead).

\begin{itemize}[leftmargin=1.4em]
  \item Leap years: 24
  \item Ordinary years: $100 - 24 = 76$
  \item Odd days $= (24 \times 2) + (76 \times 1) = 48 + 76 = 124$
  \item $124 \bmod 7 = 124 - 119 = 5$
\end{itemize}

Monday $+ 5$ odd days $\to$ \textbf{Saturday}.
\textbf{Answer: 1 January 2101 is a Saturday.}
\end{tipbox}

\begin{warnbox}[Why This Is a Common CAT/Placement Trap]
The trap is treating 2100 as a normal leap year because it's divisible by 4. Since it's a century year
\textbf{not} divisible by 400, it must be excluded --- this single oversight shifts the final answer
by one full odd day (which changes the answer from Saturday to Sunday if missed).
\end{warnbox}

\section*{Answer Key Summary}
\begin{center}
\begin{tabular}{cl}
\toprule
Question & Answer \\
\midrule
Q1 & Thursday \\
Q2 & Tuesday \\
Q3 & Tuesday \\
Q4 & 2034 \\
Q5 & Only Sun/Mon/Wed/Fri possible for century-end (proof-based) \\
Q6 & Certain year is leap; 2 years later = Monday \\
Q7 & Sunday \\
Q8 & Saturday \\
\bottomrule
\end{tabular}
\end{center}

\end{document}
